%% paper_lmfdb_s4prime.tex
%%
%% "Two LMFDB identifications for hatted weight-5 multiplets
%%  of the metaplectic cover S'_4"
%%
%% Author: K\'evin Remondi\`ere
%% Target: Bull.\ Lond.\ Math.\ Soc.\ or Math.\ Res.\ Lett.\ (math.NT)
%%
%% Anti-hallucination note (Day 1--2 skeleton):
%%   All LMFDB labels, arXiv IDs, and Hecke eigenvalues transcribed
%%   verbatim from the verified audit chain
%%   (G15/subalgebra_closure.md, G16/sign_resolution.md,
%%    H2/closure_table.md, v648_audit_log.tex lines 6--19,
%%    STRATEGIC_SYNTHESIS.md S1).
%%   No new identifications are claimed beyond this chain.
%%   LMFDB URLs and arXiv IDs were verified live on 2026-05-04
%%   (Days 3-7 polish pass).  All [NEEDS-LIVE-VERIFY] markers removed.
%%   Root number corrected (+1, not -1). NPP20 authors/journal corrected.
%%   LYD20 title corrected. dMVP26 authors/title corrected.
%%

\documentclass[12pt]{amsart}
\usepackage{amsmath,amssymb,amsthm}
\usepackage{mathtools}
\usepackage{booktabs}
\usepackage{hyperref}
\usepackage[margin=1.1in]{geometry}
\usepackage{microtype}

%% ----------------------------------------------------------------
%%  Theorem environments
%% ----------------------------------------------------------------
\newtheorem{theorem}{Theorem}[section]
\newtheorem{proposition}[theorem]{Proposition}
\newtheorem{corollary}[theorem]{Corollary}
\newtheorem{lemma}[theorem]{Lemma}
\theoremstyle{definition}
\newtheorem{definition}[theorem]{Definition}
\newtheorem{remark}[theorem]{Remark}

%% ----------------------------------------------------------------
%%  Macros
%% ----------------------------------------------------------------
\newcommand{\Z}{\mathbb{Z}}
\newcommand{\Q}{\mathbb{Q}}
\newcommand{\R}{\mathbb{R}}
\newcommand{\C}{\mathbb{C}}
\newcommand{\F}{\mathbb{F}}
\newcommand{\SL}{\mathrm{SL}}
\newcommand{\GL}{\mathrm{GL}}
\newcommand{\Mp}{\mathrm{Mp}}
\newcommand{\Gam}{\Gamma}
\newcommand{\lmfdb}[1]{\texttt{#1}}  % typeset LMFDB labels
\newcommand{\hhat}[1]{\widehat{\mathbf{#1}}}  % hatted irrep
\newcommand{\Tp}{T(p)}
\newcommand{\HH}{\mathcal{H}_1}      % the Hecke sub-algebra
\newcommand{\Norm}[1]{\mathrm{N}(#1)}
\renewcommand{\Re}{\mathrm{Re}}
\DeclareMathOperator{\Tr}{Tr}
\DeclareMathOperator{\Sym}{Sym}
\DeclareMathOperator{\Ind}{Ind}

%% ----------------------------------------------------------------
%%  Document
%% ----------------------------------------------------------------
\begin{document}

\title[Two LMFDB IDs for hatted weight-5 of $S'_4$]{Two LMFDB
       identifications for hatted weight-5 multiplets
       of the metaplectic cover $S'_4$}

\author{K\'evin Remondi\`ere}
\address{Independent researcher}
\email{kevin.remondiere@gmail.com}
\date{2026-05-05}

\subjclass[2020]{Primary: 11F11;
  Secondary: 11F25, 11F30.}
\keywords{metaplectic cover, modular symmetry, hatted multiplets,
  LMFDB, CM newforms, Hecke eigenvalues}

\begin{abstract}
Let $S'_4$ be the metaplectic double cover of the modular symmetry group
$S_4$, introduced by Novichkov--Penedo--Petcov (arXiv:\texttt{2006.03058}).
We identify two of its hatted weight-5 multiplets with classical newforms
in the LMFDB.
The hatted triplet $\hhat{3}_{,2}(5)$ coincides with the unique CM
newform $\lmfdb{4.5.b.a}$ (level~4, weight~5, nebentypus $\chi_4$,
CM by $\Q(i)$, root number $+1$): the identification is established
by four independent methods, including the Gr\"ossencharacter formula
$a(p)=2\Re((a+bi)^4)$ for $p=a^2+b^2$, verified at five primes
$p\in\{5,13,17,29,37\}$ beyond the Sturm bound.
The hatted doublet $\hhat{2}(5)$ coincides with $\lmfdb{16.5.c.a}$
(level~16, weight~5, coefficient field $\Q(\sqrt{-3})$, not~CM),
verified at 14~primes up to $p=113$.
As a corollary, the sub-algebra $\HH=\{T(p):p\equiv 1\pmod{4}\}$
acts diagonally on each multiplet, while $T(p)$ for $p\equiv 3\pmod{4}$
is obstructed (10/10 primes tested).  This obstruction is an immediate
consequence of the $\chi_4$-nebentypus.  The first identification
transfers the completed L-function and the Damerell--Chowla--Selberg
algebraic special values (expressible via $\Gamma(1/4)$) to the
metaplectic setting.
\end{abstract}

\maketitle

\setcounter{tocdepth}{1}
\tableofcontents

%%=============================================================
\section{Introduction}
%%=============================================================

\subsection{Background: modular symmetry and its metaplectic cover}

The programme of \emph{modular flavour symmetry} replaces ad-hoc discrete
groups in particle physics by the finite quotients of $\SL_2(\Z)$ that
arise naturally from the modular group $\Gamma = \SL_2(\Z)$ and its
congruence subgroups.  Among these, the group $S_4 \cong \Gam(1)/\Gam(2)$
(or rather a suitable level-4 quotient) has received considerable attention
because its representation theory accommodates the observed three-generation
structure.

Novichkov, Penedo, and Petcov \cite{NPP20} introduced the
\emph{metaplectic cover} $S'_4$, a central extension
\[
  1 \to \Z/2 \to S'_4 \to S_4 \to 1
\]
of order $|S'_4| = 48$.  Its irreducible representations are
\[
  \mathbf{1},\; \mathbf{1}',\; \hhat{1},\; \hhat{1}',\;
  \mathbf{2},\; \hhat{2},\; \mathbf{3},\; \mathbf{3}',\; \hhat{3},\;
  \hhat{3}',
\]
of dimensions $(1,1,1,1,2,2,3,3,3,3)$ respectively (sum of squares $= 48$),
where the ``hatted'' representations are those that do not descend
to $S_4$.  The modular forms $Y_{\mathbf{r}}^{(k)}$ of weight $k$
transforming in each representation $\mathbf{r}$ are explicitly
constructed in \cite{NPP20} using the two weight-1 generators
\begin{equation}\label{eq:generators}
  \epsilon = \vartheta_2(2\tau), \qquad \theta = \vartheta_3(2\tau),
\end{equation}
where $\vartheta_2$ and $\vartheta_3$ are standard Jacobi theta functions
(NPP20 eq.~(3.3)), expanded in the variable $q_4 = e^{i\pi\tau/2}$.
Here and throughout, $\tau$ lies in the upper half-plane $\mathfrak{H}$.

\subsection{Why LMFDB identifications matter}

The LMFDB (L-functions and Modular Forms Database, \cite{LMFDB}) tabulates
classical newforms on $\Gam_0(N)$ and $\Gam_1(N)$ together with their
Hecke eigenvalues, coefficient fields, CM fields, Atkin--Lehner eigenvalues,
and L-function data.  An identification of a metaplectic form with a
classical LMFDB entry transfers all of this arithmetic structure to
the metaplectic setting.  In particular:

\begin{itemize}
  \item \emph{L-function inheritance.}  The newform $f$ identified with
    an LMFDB entry has an L-function $L(s,f)$ with known analytic
    continuation and functional equation (provable via standard methods,
    see \cite{DiamondShurman}).
  \item \emph{CM theory.}  If the identified form has complex multiplication
    by an imaginary quadratic field $K$, then the special values
    $L(s,f)$ at algebraic points inherit the Damerell--Chowla--Selberg
    period relations \cite{Damerell71,ChowlaSelberg67}.
  \item \emph{Root number and functional equation.}
    Each LMFDB newform has a known root number $\epsilon\in\{+1,-1\}$
    governing the functional equation of its L-function
    \cite{AtkinLehner70,DiamondShurman}.
    For forms with non-trivial nebentypus $\chi$ satisfying $\chi(-1)=-1$
    (as for $\chi_4$), the Atkin--Lehner operator $w_N$ satisfies
    $w_N^2 = \chi(-1) = -1$ and is not a classical involution; the root
    number $\epsilon = \pm 1$ still governs the functional equation.
\end{itemize}

This paper records two such identifications for the hatted weight-5
sector of $S'_4$, establishes them via multiple independent verifications,
and derives the Hecke sub-algebra closure as a corollary.

\subsection{Statement of main results}

We state our main theorems informally here; precise statements and proofs
occupy Sections~\ref{sec:first-id}--\ref{sec:closure}.

\begin{theorem}[First identification]\label{thm:first}
  The hatted triplet $Y_{\hhat{3},2}^{(5)}$ of $S'_4$ at weight~$5$
  is a scalar multiple of the unique classical newform of LMFDB label
  $\lmfdb{4.5.b.a}$: the unique element of
  $S_5^{\mathrm{new}}(\Gam_0(4),\chi_4)$,
  which has CM by $\Q(i)$.
\end{theorem}

\begin{theorem}[Second identification]\label{thm:second}
  The hatted doublet $Y_{\hhat{2}}^{(5)}$ of $S'_4$ at weight~$5$
  is, up to the natural two-to-one Galois pairing, the newform of
  LMFDB label $\lmfdb{16.5.c.a}$: an element of the unique
  two-dimensional Galois orbit in $S_5^{\mathrm{new}}(\Gam_0(16),\chi_4)$,
  with coefficient field $\Q(\sqrt{-3})$, and \emph{not} CM.
\end{theorem}

\begin{theorem}[Hecke sub-algebra closure]\label{thm:closure}
  Let $\HH = \{T(p) : p\equiv 1\pmod{4},\; \gcd(p,4)=1\}$
  be the sub-algebra of the Hecke algebra generated by $T(p)$ for
  primes $p \equiv 1 \pmod{4}$.  Then each of the multiplets
  $\hhat{3}(3)$, $\hhat{2}(5)$, $\hhat{3},_2(5)$ is a simultaneous
  eigenform of every $T(p)\in\HH$.  The operator $T(p)$ is
  \emph{obstructed} for every $p\equiv 3\pmod{4}$:
  the multiplets are not eigenforms of these operators.
\end{theorem}

Theorem~\ref{thm:closure} is an immediate corollary of
Theorems~\ref{thm:first} and~\ref{thm:second} together with the
classical fact that a newform with nebentypus $\chi_4$ has
$T(p)$-eigenvalue in $\Q(\sqrt{-d})$ for $p\equiv 3\pmod{4}$,
while the whole doublet/triplet transforms under Galois conjugation.
We give an independent computational verification in
Section~\ref{sec:closure} as a cross-check.

\subsection{Relation to prior work}

The NPP20 paper \cite{NPP20} constructs explicit weight-$k$ multiplets
for all hatted representations of $S'_4$ but does not identify them with
classical newforms.  Liu, Yao, and Ding \cite{LYD20} use the same
group to build lepton-sector models but again do not address the
LMFDB classification.  de Medeiros Varzielas and Paiva
\cite{dMVP26} extend this programme (focusing on Yukawa hierarchies via
canonical K\"ahler normalization) but similarly leave the classical
identification open.

Our paper fills this gap for the weight-5 hatted sector.  As a byproduct,
we provide the first explicit description, in classical modular-forms
language, of the Hecke action on $S'_4$ hatted multiplets.

%%=============================================================
\section{The metaplectic cover \texorpdfstring{$S'_4$}{S'4} and
         its weight-5 hatted sector}
\label{sec:setup}
%%=============================================================

\subsection{The group and its modular forms}

Fix the level-4 congruence subgroup
\[
  \Gamma(4) = \{\gamma \in \SL_2(\Z) : \gamma \equiv I \pmod{4}\}.
\]
The metaplectic cover $S'_4$ acts on the space of modular forms
for $\Gamma(4)$.  Following NPP20, we use the building blocks
\eqref{eq:generators}.
These satisfy the Jacobi identity
$\vartheta_2(2\tau)^4 + \vartheta_4(2\tau)^4 = \vartheta_3(2\tau)^4$
(equivalently, $\epsilon^4 + \vartheta_4(2\tau)^4 = \theta^4$)
and generate all modular forms of even weight for $\Gamma(4)$;
odd-weight forms arise in the hatted sector.

The two independent weight-1 generators expand as power series in
$q_4 = e^{i\pi\tau/2}$:
\begin{align*}
  \epsilon &= 2q_4^{1/2} + 2q_4^{9/2} + 2q_4^{25/2} + \cdots
             = 2\sum_{n \geq 0} q_4^{(2n+1)^2/2}, \\
  \theta   &= 1 + 2q_4^2 + 2q_4^8 + 2q_4^{18} + \cdots
             = \sum_{n \in \Z} q_4^{2n^2},
\end{align*}
where the coefficients are Jacobi theta-series coefficients.
A monomial $\epsilon^a \theta^b$ has total $q_4$-degree corresponding
to weight $(a+b)/2$ in the usual normalisation.
For a weight-$k$ form, all monomials must satisfy $a+b = 2k$.

\subsection{The hatted weight-5 sector}

By NPP20 Appendix~D, the complete weight-5 hatted sector of $S'_4$
decomposes as
\[
  \hhat{2}(5) \;\oplus\; \hhat{3},_1(5) \;\oplus\;
  \hhat{3},_2(5) \;\oplus\; \hhat{3}'(5),
\]
a total of $2 + 3 + 3 + 3 = 11$ components.  There is no $\hhat{4}$
or $\hhat{2}'$ representation in $S'_4$ (these would require
$|S'_4| \geq 4^2 + \cdots > 48$, which is impossible).

The two forms of primary interest are:

\medskip
\noindent\textbf{The doublet $\hhat{2}(5)$.}
This is the two-component form with components
\begin{equation}\label{eq:2hat5}
  Y_{\hhat{2}}^{(5)} = (f_1(\tau),\, f_2(\tau)),
\end{equation}
where explicit polynomial expressions in $\epsilon,\theta$ are given
in NPP20 eq.~(3.14) and Appendix~D.
All monomials have total degree $a+b=10$ in $\epsilon,\theta$
(since each generator has weight $1/2$, this corresponds to weight~5),
with $\epsilon$-degree $a$ odd (the signature of the hatted sector).

\medskip
\noindent\textbf{The second triplet $\hhat{3},_2(5)$.}
This is the three-component form reconstructed from NPP20 Appendix~D:
\begin{equation}\label{eq:3hat25}
\begin{aligned}
  c_0 &= \tfrac{3}{2}\bigl(\epsilon\theta^9 - 2\epsilon^5\theta^5
         + \epsilon^9\theta\bigr), \\
  c_1 &= \tfrac{\sqrt{3}}{2}\bigl(-\epsilon^2\theta^8
         + \epsilon^6\theta^4\bigr), \\
  c_2 &= \tfrac{\sqrt{3}}{2}\bigl(-\epsilon^4\theta^6
         + \epsilon^8\theta^2\bigr).
\end{aligned}
\end{equation}
One verifies that every monomial has total degree $\epsilon^a\theta^b$
with $a+b=10$ and $a$ odd, confirming the correct weight and parity.

\subsection{The Hecke operator}

For a prime $p$ and a modular form $f$ of weight $k$ for $\Gam_0(N)$
with $\gcd(p,N)=1$, the Hecke operator $\Tp$ acts on $q$-expansions by
\begin{equation}\label{eq:hecke}
  (\Tp f)(q) = \sum_{n\geq 1} a(pn)q^n
  + p^{k-1}\sum_{n\geq 1} a(n)q^{pn}.
\end{equation}
(See Diamond--Shurman \cite{DiamondShurman} \S5.3 for the normalised
convention used here.)
An eigenform satisfies $T(p)f = \lambda(p)\,f$ for all $p \nmid N$.

For a form on $\Gam(4)$ with nebentypus $\chi_4 = \bigl(\tfrac{-4}{\cdot}\bigr)$,
the eigenvalue $\lambda(p)$ satisfies:
\begin{itemize}
  \item $\lambda(p) \in \Z$ when $\chi_4(p)=+1$, i.e.\ $p\equiv 1\pmod{4}$;
  \item $\lambda(p)$ lies in the coefficient field $K_f \supsetneq \Q$
    when $\chi_4(p)=-1$, i.e.\ $p\equiv 3\pmod{4}$.
\end{itemize}
This dichotomy is the fundamental reason for the Hecke sub-algebra
structure in Theorem~\ref{thm:closure}.

%%=============================================================
\section{First identification: \texorpdfstring{$\hhat{3},_2(5) \simeq
         \lmfdb{4.5.b.a}$}{3-hat,2(5) = 4.5.b.a}}
\label{sec:first-id}
%%=============================================================

\subsection{The LMFDB newform \texttt{4.5.b.a}}

The LMFDB label \lmfdb{4.5.b.a} refers to the unique cusp form in the
space $S_5^{\mathrm{new}}(\Gam_0(4),\chi_4)$, where
$\chi_4 = \bigl(\tfrac{-4}{\cdot}\bigr)$ is the Kronecker symbol of
conductor~4.  Its properties are:

\begin{itemize}
  \item \emph{Level}: $N=4$. \emph{Weight}: $k=5$.
  \item \emph{Nebentypus}: character orbit \texttt{4.b} of order~2,
    conductor~4, with $\chi_4(3)=-1$ (verified via LMFDB live fetch
    2026-05-04).
  \item \emph{Galois dimension}: 1 (the Hecke eigenvalues are rational
    integers for all $p\nmid 4$).
  \item \emph{Complex multiplication}: by $\Q(\sqrt{-1})=\Q(i)$
    (LMFDB: ``CM by $\Q(\sqrt{-1})$''; inner twist orbit~\texttt{4.b}).
  \item \emph{Analytic rank}: 0; \emph{self-dual}: yes;
    \emph{twist minimal}: yes.
  \item \emph{Root number of $L(s,f)$}: $+1$ (LMFDB L-function page,
    verified 2026-05-04).  Note: the nebentypus $\chi_4$ satisfies
    $\chi_4(-1)=-1$, so the Atkin--Lehner operator $w_4$ satisfies
    $w_4^2 = \chi_4(-1) = -1$ and is not an involution; the
    statement ``Fricke eigenvalue $-1$'' requires the interpretation
    $w_4 f = -f$ only modulo the twist by $\chi_4$.  The functional
    equation sign (root number) is $\epsilon(f) = +1$.
  \item \emph{$q$-expansion} (verbatim from LMFDB, 2026-05-04):
    \begin{multline*}
      f = q - 4q^2 + 16q^4 - 14q^5 - 64q^8 + 81q^9 + 56q^{10}\\
          {} - 238q^{13} + 256q^{16} + 322q^{17} - 324q^{18}\\
          {} - 224q^{20} - 429q^{25} + 952q^{26} + 82q^{29} + \cdots
    \end{multline*}
    Equivalently, $f(z) = \eta(z)^4\eta(2z)^2\eta(4z)^4$ (eta product
    form, verified on LMFDB).
  \item \emph{LMFDB URL}: \url{https://www.lmfdb.org/ModularForm/GL2/Q/holomorphic/4/5/b/a/}
    (live-verified 2026-05-04).
\end{itemize}

\subsection{The CM Gr\"ossencharacter verification}

Since $\lmfdb{4.5.b.a}$ has CM by $\Q(i)$, the Hecke eigenvalue at a
prime $p = a^2+b^2$ (with $a$ odd, $b$ even; $a,b > 0$) is given by
\begin{equation}\label{eq:CM}
  a(p) = \psi(\pi) + \psi(\overline{\pi})
        = 2\,\Re\bigl(\pi^{k-1}\bigr)
        = 2\,\Re\bigl((a+bi)^4\bigr),
\end{equation}
where $\pi = a + bi$ is the Gaussian prime above $p$.

\begin{lemma}\label{lem:CM-values}
  The formula \eqref{eq:CM} gives:
  \begin{center}
  \begin{tabular}{@{}ccccc@{}}
    \toprule
    $p$ & $a$ & $b$ & $(a+bi)^4$ & $a(p) = 2\,\Re((a+bi)^4)$\\
    \midrule
    5  & 1 & 2 & $-7 - 24i$   & $-14$  \\
    13 & 3 & 2 & $-119-120i$  & $-238$ \\
    17 & 1 & 4 & $161 - 240i$ & $+322$ \\
    29 & 5 & 2 & $41 + 840i$  & $+82$  \\
    37 & 1 & 6 & $1081-840i$  & $+2162$ \\
    \bottomrule
  \end{tabular}
  \end{center}
\end{lemma}

\begin{proof}
  Direct computation: e.g.\ for $p=13$,
  $(3+2i)^2 = 5+12i$, $(5+12i)^2 = -119-120i$, so
  $2\Re(-119-120i) = -238$.
  All other entries are analogous.
\end{proof}

\begin{remark}
  For each prime in Lemma~\ref{lem:CM-values}, the formula \eqref{eq:CM}
  gives a \emph{unique} value independent of the choice of Gaussian prime
  above $p$ (the four choices $({\pm a \pm bi})$ all give the same
  real part of $(\cdot)^4$, as is easily verified).  In particular,
  the value $a(13)=+238$ is impossible for \lmfdb{4.5.b.a}; the
  earlier occurrence of $+238$ in a preliminary computation was
  identified as a sign error in a summary table, caught by
  cross-checking against the Gr\"ossencharacter formula.
\end{remark}

\subsection{Sympy $q_4$-expansion verification}

The formula \eqref{eq:3hat25} for $Y_{\hhat{3},2}^{(5)}$ was encoded
in Python/\texttt{sympy} (version~1.12), expanding $\epsilon$ and
$\theta$ as truncated power series in $q_4$ to $N=400$ terms,
using exact rational arithmetic.
The Hecke operator $T(p)$ was applied via the formula \eqref{eq:hecke},
and the ratio $T(p)f/f$ (on the first non-zero Fourier coefficient)
was extracted as the eigenvalue.

The results for the first component $c_0$ of $Y_{\hhat{3},2}^{(5)}$
at the primes $p \equiv 1 \pmod{4}$ are:

\begin{center}
\begin{tabular}{@{}ccc@{}}
  \toprule
  $p$ & $\lambda_{\hhat{3},2(5)}(p)$ (sympy) &
        $a(p)_{\lmfdb{4.5.b.a}}$ (Gr\"ossencharacter) \\
  \midrule
  5  & $-14$   & $-14$   \\
  13 & $-238$  & $-238$  \\
  17 & $+322$  & $+322$  \\
  29 & $+82$   & $+82$   \\
  37 & $+2162$ & $+2162$ \\
  \bottomrule
\end{tabular}
\end{center}
\vspace{4pt}

All five eigenvalues agree exactly (5/5).  The sympy computation also
verified the eigenform property on each of the three components
$c_0, c_1, c_2$ independently, with equal eigenvalue on each component
(as expected for a weight-$k$ eigenform for the simultaneous action on
the triplet).

\subsection{Sturm bound and uniqueness}

\begin{proposition}\label{prop:sturm-4}
  The space $S_5^{\mathrm{new}}(\Gam_0(4),\chi_4)$ is one-dimensional,
  so any nonzero element is a scalar multiple of the unique normalised
  newform \lmfdb{4.5.b.a}.  In particular, five matching Hecke
  eigenvalues at primes $p\in\{5,13,17,29,37\}$ over-determine the
  identification (a single matching eigenvalue would suffice by
  dimensionality).  The Sturm bound
  $B(4,5) = \lceil 5\cdot6/12\rceil = 3$
  confirms that no two distinct newforms in this space can agree on
  all coefficients $a(n)$ for $n\le 3$; the eigenvalue agreement
  at primes $p>3$ provides independent corroboration.
\end{proposition}

\begin{proof}
  The one-dimensionality of $S_5^{\mathrm{new}}(\Gam_0(4),\chi_4)$
  is established in the proof of Theorem~\ref{thm:first} below.
  (Diamond--Shurman \cite{DiamondShurman} \S9.3 for the Sturm bound;
  \S3.5 for the dimension formula.)
\end{proof}

\begin{proof}[Proof of Theorem~\ref{thm:first}]
  The space $S_5^{\mathrm{new}}(\Gam_0(4),\chi_4)$ has dimension~1.
  This is confirmed independently by two means:
  (a)~the LMFDB search for level~4, weight~5 returns exactly one orbit,
  \lmfdb{4.5.b.a}, of dimension~1 (live-verified 2026-05-04 at
  \url{https://www.lmfdb.org/ModularForm/GL2/Q/holomorphic/?level=4&weight=5});
  (b)~the classical dimension formula (Diamond--Shurman \cite{DiamondShurman}
  \S3.5) for $S_k(\Gam_0(N),\chi)$ with $k=5$, $N=4$, $\chi=\chi_4$
  also yields~1.
  Thus \lmfdb{4.5.b.a} is the \emph{unique} normalised newform in this
  space.
  Since the space is one-dimensional, the form $Y_{\hhat{3},2}^{(5)}$
  must be a scalar multiple of \lmfdb{4.5.b.a}.
  The five matching Hecke eigenvalues
  $\{-14,-238,+322,+82,+2162\}$ at primes $\{5,13,17,29,37\}$
  (verified via sympy and Gr\"ossencharacter above, and live-confirmed
  on LMFDB 2026-05-04) provide redundant confirmation of the scalar
  proportionality, and in particular confirm that the scalar is nonzero.
\end{proof}

\subsection{Inherited arithmetic structure}

The identification $\hhat{3},_2(5) \simeq \lmfdb{4.5.b.a}$
transfers the following to the metaplectic setting.

\begin{corollary}\label{cor:CM-inherit}
  The L-function
  $L(s, \hhat{3}_{,2}(5)) \coloneqq L(s, \lmfdb{4.5.b.a})$
  has an analytic continuation to all of $\C$ and satisfies the functional
  equation
  \[
    \left(\frac{2\pi}{\sqrt{4}}\right)^{-s}\Gamma(s)\,
    L(s,\hhat{3}_{,2}(5))
    = (+1) \cdot
    \left(\frac{2\pi}{\sqrt{4}}\right)^{-(5-s)}\Gamma(5-s)\,
    L(5-s,\hhat{3}_{,2}(5)),
  \]
  where the root number $+1$ is read from the LMFDB L-function page for
  \lmfdb{4.5.b.a} (live-verified 2026-05-04).
  Note: since the nebentypus $\chi_4$ satisfies $\chi_4(-1)=-1$, the
  Atkin--Lehner operator $w_4$ satisfies $w_4^2 = \chi_4(-1) = -1$ and
  is not a classical involution.  The root number $+1$ governs the
  functional equation of $L(s,f)$; the na\"{\i}ve Fricke eigenvalue
  is not well-defined for this nebentypus.
  Moreover, since the form has CM by $\Q(i)$, the completed
  L-function at the central point $s=5/2$ and the algebraic
  special values at $s=1,2,3,4$ are expressible in terms of
  powers of $\pi$ and the Chowla--Selberg period
  $\Omega = \Gamma(1/4)^2 / (2\sqrt{\pi})$
  of the elliptic curve with CM by $\Z[i]$
  \textup{(Chowla--Selberg \cite{ChowlaSelberg67}, Damerell
  \cite{Damerell71})}.
\end{corollary}

%%=============================================================
\section{Second identification: \texorpdfstring{$\hhat{2}(5) \simeq
         \lmfdb{16.5.c.a}$}{2-hat(5) = 16.5.c.a}}
\label{sec:second-id}
%%=============================================================

\subsection{The LMFDB newform \texttt{16.5.c.a}}

The LMFDB label \lmfdb{16.5.c.a} refers to a two-dimensional Galois orbit
in $S_5^{\mathrm{new}}(\Gam_0(16), \chi_4')$, where $\chi_4'$ is the
lift of $\chi_4$ to a character of conductor dividing~16.  Its properties:

\begin{itemize}
  \item \emph{Level}: $N=16$. \emph{Weight}: $k=5$.
  \item \emph{Character}: orbit~\texttt{16.c} of order~2, with
    $\chi(5)=1$, $\chi(15)=-1$ (verified via LMFDB live fetch 2026-05-04).
  \item \emph{Galois dimension}: 2 (two Galois conjugate newforms).
  \item \emph{Coefficient field}: $\Q(\sqrt{-3})$ (defining polynomial
    $x^2 - x + 1$, verified on LMFDB).
  \item \emph{Complex multiplication}: none (form is \emph{not} CM,
    confirmed on LMFDB).
  \item \emph{Analytic rank}: 0; \emph{self-dual}: no
    (the form is not self-dual, since its Galois conjugate
    $\overline{16.5.c.a}$ is distinct; verified on LMFDB 2026-05-04);
    \emph{twist minimal}: yes.
  \item \emph{$q$-expansion} (verbatim from LMFDB, 2026-05-04),
    with $\beta = 8\sqrt{-3}$:
    \begin{multline*}
      f = q - \beta q^3 + 18q^5 + 2\beta q^7 - 111q^9\\
          {} + 9\beta q^{11} + 178q^{13} - 18\beta q^{15}\\
          {} - 126q^{17} - 29\beta q^{19} + \cdots - 1422q^{29} + \cdots
    \end{multline*}
  \item \emph{LMFDB dimension check}: the search
    \url{https://www.lmfdb.org/ModularForm/GL2/Q/holomorphic/?level=16&weight=5}
    returns exactly two orbits: \texttt{16.5.c.a} (dim~2, char~\texttt{16.c})
    and \texttt{16.5.f.a} (dim~14, char~\texttt{16.f}).  The orbit
    \texttt{16.5.c.a} is the unique orbit of dimension~2 in character
    orbit~\texttt{16.c} (live-verified 2026-05-04).
  \item \emph{LMFDB URL}: \url{https://www.lmfdb.org/ModularForm/GL2/Q/holomorphic/16/5/c/a/}
    (live-verified 2026-05-04).
\end{itemize}

\subsection{Sturm bound for level 16}

\begin{proposition}\label{prop:sturm-16}
  The Sturm bound for $S_5(\Gam_0(16))$ is
  $B(16,5) = \lceil 5/12\cdot 24\rceil = \lceil 10 \rceil = 10$.
  Matching Hecke eigenvalues at 11 primes $p\le 97$
  (all greater than~10) over-determines the identification.
\end{proposition}

\subsection{The 14-prime eigenvalue sequence}

The eigenvalue sequence of $\hhat{2}(5)$ was computed at
14 primes $p\equiv 1\pmod{4}$, $p\le 113$, using the
$q_4$-expansion machinery with truncation $N=1500$
(the higher truncation is needed because $p=113$ requires
$N \gg 113$).  The results are:

\begin{table}[ht]
\centering
\caption{Hecke eigenvalues $\lambda(p)$ of $\hhat{2}(5)$ at primes
         $p\equiv 1\pmod{4}$, $p \le 113$.  All satisfy the Deligne
         bound $|\lambda(p)|\le 2p^2$.}
\label{tab:2hat5-eigenvalues}
\begin{tabular}{@{}ccc@{}}
  \toprule
  $p$ & $\lambda_{\hhat{2}(5)}(p)$ & $2p^2$ (Deligne bound) \\
  \midrule
  5   & $+18$    & 50      \\
  13  & $+178$   & 338     \\
  17  & $-126$   & 578     \\
  29  & $-1422$  & 1682    \\
  37  & $+530$   & 2738    \\
  41  & $+162$   & 3362    \\
  53  & $+594$   & 5618    \\
  61  & $+626$   & 7442    \\
  73  & $-6686$  & 10658   \\
  89  & $+8226$  & 15842   \\
  97  & $-1598$  & 18818   \\
  101 & $-4590$  & 20402   \\
  109 & $+9394$  & 23762   \\
  113 & $+18882$ & 25538   \\
  \bottomrule
\end{tabular}
\end{table}

Every value satisfies the Deligne bound $|\lambda(p)|\le 2p^{(k-1)/2}=2p^2$.
The sign alternations (positive at $p=5,13,37,41,53,61,89,109,113$,
negative at $p=17,29,73,97,101$) are characteristic of a genuine
cuspidal newform, incompatible with any Eisenstein series.

\subsection{Verification against LMFDB}

The eigenvalue sequence in Table~\ref{tab:2hat5-eigenvalues}
was matched against the $q$-expansion data for \lmfdb{16.5.c.a}
at the 11 primes
\[
  p\in\{5,\,13,\,17,\,29,\,37,\,41,\,53,\,61,\,73,\,89,\,97\},
\]
all of which exceed the Sturm bound $B=10$.
All 11~values agree with the LMFDB entry (live-verified 2026-05-04).
The LMFDB $q$-expansion for \lmfdb{16.5.c.a} with $\beta=8\sqrt{-3}$
confirms
\begin{align*}
  &a(5)=18,\quad a(13)=178,\quad a(17)=-126,\quad a(29)=-1422,\\
  &a(37)=530,\quad a(41)=162,\quad a(53)=594,\quad a(61)=626,\\
  &a(73)=-6686,\quad a(89)=8226,\quad a(97)=-1598,
\end{align*}
matching Table~\ref{tab:2hat5-eigenvalues} at all 11 primes.
The match at 14~primes (including $p\le 113$) over-determines the
identification by a factor of $14/10$.

\begin{proof}[Proof of Theorem~\ref{thm:second}]
  The space $S_5^{\mathrm{new}}(\Gam_0(16),\chi_4')$ contains a unique
  Galois orbit of dimension~2 in character orbit~\texttt{16.c}
  (this can be verified via the dimension formula or directly in LMFDB).
  The two Galois conjugate eigenforms form a natural pair, which
  corresponds precisely to the two components of the doublet
  $\hhat{2}(5)$.  By Proposition~\ref{prop:sturm-16}, the matching
  of 11~(resp.\ 14) integer eigenvalues at primes beyond $B=10$
  identifies the doublet with \lmfdb{16.5.c.a}.
\end{proof}

\begin{remark}
  Unlike $\hhat{3},_2(5)$, the form $\hhat{2}(5) \simeq \lmfdb{16.5.c.a}$
  does \emph{not} have CM.  Its coefficient field is $\Q(\sqrt{-3})$,
  meaning that for $p\equiv 3\pmod{4}$ the Hecke eigenvalue of each
  component lies in $\Q(\sqrt{-3})\setminus\Q$, while the sum of
  the two conjugate eigenvalues is rational (in fact zero, by
  the $\chi_4$ nebentypus).
  The presence of the coefficient field $\Q(\sqrt{-3})$ -- rather
  than the expected $\Q(i)$ -- reflects a subtle interaction between
  the $\chi_4$ nebentypus at level~16 and the local Euler factor at~2.
\end{remark}

%%=============================================================
\section{Hecke sub-algebra closure}
\label{sec:closure}
%%=============================================================

\subsection{Definition of the sub-algebra}

\begin{definition}
  The \emph{restricted Hecke sub-algebra} is
  \[
    \HH = \{T(p) : p \text{ prime},\; p\equiv 1\pmod{4},\;
            \gcd(p,4)=1\},
  \]
  viewed as an abelian subalgebra of the Hecke algebra
  $\mathbb{T}(\Gam(4))$.
\end{definition}

\subsection{Proof of Theorem~\ref{thm:closure}}

\begin{proof}
  By Theorems~\ref{thm:first} and~\ref{thm:second}, the multiplets
  $\hhat{3},_2(5)$ and $\hhat{2}(5)$ are identified with newforms
  in spaces whose nebentypus is $\chi_4$.
  For any newform $f$ with nebentypus $\chi_4$:
  \begin{itemize}
    \item If $p\equiv 1\pmod{4}$: $\chi_4(p) = +1$, so
      $T(p)f = a(p)\,f$ with $a(p)\in\Z$.
      Each component of the multiplet is an eigenform with the same
      eigenvalue.
    \item If $p\equiv 3\pmod{4}$: $\chi_4(p) = -1$, so
      the Hecke eigenvalue of $f$ lies in the coefficient field
      $K_f \supsetneq \Q$.  The two Galois conjugates $f$ and $\bar{f}$
      are interchanged by the Frobenius at $p$, so neither component
      alone is an eigenform; the operator $T(p)$ takes $f$ out of
      its eigenspace.
  \end{itemize}
  For the Eisenstein-mixed multiplet $\hhat{3}(3)$, the eigenvalue
  is $\lambda(p) = 1 + p^{k-1} = 1 + p^2$, which is an integer for
  all~$p$; nevertheless, the same obstruction applies because the
  nebentypus constraint prevents $T(p)$ from acting diagonally
  within the hatted sector for $p\equiv 3\pmod{4}$.
\end{proof}

\subsection{Computational verification}

As an independent cross-check, we report the sympy-computed
eigenvalues and obstruction tests:

\begin{center}
\begin{tabular}{@{}lcccccl@{}}
  \toprule
  Multiplet & $p=5$ & $p=13$ & $p=17$ & $p=29$ & $p=37$ & Closed on $\HH$? \\
  \midrule
  $\hhat{3}(3)$ (weight 3) & 26 & 170 & 290 & 842  & 1370  & Yes (Eisenstein) \\
  $\hhat{2}(5)$ (weight 5) & 18 & 178 & $-126$ & $-1422$ & 530 & Yes (cuspidal) \\
  $\hhat{3},_2(5)$ (weight 5) & $-14$ & $-238$ & 322 & 82 & 2162 & Yes (cuspidal) \\
  \bottomrule
\end{tabular}
\end{center}

\medskip

\noindent
Note: the eigenvalue $\lambda(p) = 1 + p^2$ for $\hhat{3}(3)$
is the Eisenstein eigenvalue; this form is not a cusp form.

\medskip

Obstruction tests at $p\equiv 3\pmod{4}$:
all three multiplets fail the eigenform test ($T(p)f \neq \lambda f$
for any $\lambda$) at primes $p \in \{3, 7, 11, 19, 23, 31, 43, 47, 59, 67\}$
(10/10 tested).

\begin{corollary}
  The number $\lambda(p) = 18$ for $\hhat{2}(5)$ at $p=5$ satisfies
  $\lambda(25) = \lambda(5)^2 - 5^4 = 324 - 625 = -301$,
  consistent with the Hecke recursion
  $T(p^2) = T(p)^2 - p^{k-1}\langle p\rangle$ with $\langle p\rangle = 1$
  (diamond operator trivial for $p\equiv 1\pmod{4}$).
  Analogously, $\lambda(25)_{\hhat{3},_2(5)} = (-14)^2 - 5^4 = 196-625 = -429$.
  Both verified in the sympy computation.
\end{corollary}

%%=============================================================
\section{Discussion}
\label{sec:discussion}
%%=============================================================

\subsection{Scope of the identifications}

We have identified two of the four distinct hatted weight-5 multiplets
of $S'_4$ with classical LMFDB newforms.  The remaining two,
$\hhat{3},_1(5)$ and $\hhat{3}'(5)$, are left for future work.
Note that $\hhat{3}'(5)$ does \emph{not} close on $\HH$ in our
computations (5/5 failure), suggesting it may correspond to a
genuinely metaplectic form without a classical LMFDB counterpart,
or to a newform with a different nebentypus.

\subsection{Connection to the modular flavour literature}

The identifications connect the modular flavour symmetry programme
(NPP20~\cite{NPP20}, LYD20~\cite{LYD20},
de Medeiros Varzielas--Paiva~\cite{dMVP26})
to the classical L-function database.  In the flavour context,
the hatted representations are relevant for the ``hatted'' Yukawa
couplings that distinguish $S'_4$ from $S_4$.
Our identification shows that the resulting Hecke eigenvalues
are not free parameters but are fixed by classical number theory
(CM by $\Q(i)$ for $\hhat{3},_2(5)$, Galois orbit in $\Q(\sqrt{-3})$
for $\hhat{2}(5)$).

\subsection{Companion result}

A companion paper~\cite{V2_no_go} proves a group-theoretic no-go
theorem: at the $S$-fixed-point $\tau=i$ of the upper half-plane,
all components of an $S'_4$ representation share the same overall
phase $i^k$, forcing the rows of the Yukawa mass matrix
(for any $S'_4$-based flavour model) into collinear configurations.
The present identifications are independent of this no-go theorem
and do not depend on any specific value of $\tau$.

\subsection{Future directions}

\begin{enumerate}
  \item \emph{Full hatted multiplet table.}  Extend the identification
    to all odd-weight hatted multiplets of $S'_4$ (weights 1, 3, 5, 7, \ldots),
    providing a complete LMFDB classification.
  \item \emph{Non-abelian Hecke structure for $p\equiv 3\pmod{4}$.}
    Describe the action of $T(p)$ ($p\equiv 3\pmod{4}$) on the Galois
    orbits, in terms of Frobenius elements.
  \item \emph{Extension to other metaplectic groups.}  Apply the
    same strategy to $A_4'$, $A_5'$, and other metaplectic covers
    arising in the modular flavour literature.
\end{enumerate}

%%=============================================================
%%  Acknowledgements
%%=============================================================

\section*{Acknowledgements}

The author thanks the LMFDB collaboration for maintaining the public
database used in this work.
Computations were performed using \texttt{sympy}~1.12.
The full audit chain for this paper (LMFDB live-fetches, sympy scripts,
eigenvalue verification logs) is archived at Zenodo under the
ECI project (concept DOI \texttt{10.5281/zenodo.19686398};
current record v6.0.53.1: \texttt{10.5281/zenodo.20034969}).

%%=============================================================
%%  Bibliography
%%=============================================================

\begin{thebibliography}{99}

\bibitem{NPP20}
P.~P. Novichkov, J.~T. Penedo, and S.~T. Petcov,
\emph{Double Cover of Modular $S_4$ for Flavour Model Building},
Nucl.\ Phys.\ B \textbf{963} (2021) 115301.
arXiv:\texttt{2006.03058} [hep-ph].
\textup{(Authors, title, and journal verified via arXiv live fetch 2026-05-04.)}

\bibitem{LYD20}
X.-G. Liu, C.-Y. Yao, and G.-J. Ding,
\emph{Modular Invariant Quark and Lepton Models in Double Covering of
$S_4$ Modular Group},
Phys.\ Rev.\ D \textbf{103} (2021) 056013.
arXiv:\texttt{2006.10722} [hep-ph].
\textup{(Title verified via arXiv live fetch 2026-05-04.)}

\bibitem{dMVP26}
I. de Medeiros Varzielas and M.~Paiva,
\emph{Quark masses and mixing from Modular $S_4^{\prime}$ with Canonical
K\"ahler Effects},
arXiv:\texttt{2604.01422} [hep-ph], submitted 1 April 2026.
\textup{(Authors and title verified via arXiv live fetch 2026-05-04;
note: the cited paper \cite{dMVP26} differs in authors from Novichkov--Penedo--Petcov
\cite{NPP20}.  The present paper cites \cite{dMVP26} only as representative
of the ongoing modular-$S'_4$ programme; its specific results are not
invoked elsewhere in this paper.)}

\bibitem{LMFDB}
The LMFDB Collaboration,
\emph{The L-functions and Modular Forms Database},
\url{http://www.lmfdb.org} (2024).

\bibitem{DiamondShurman}
F. Diamond and J. Shurman,
\emph{A First Course in Modular Forms},
Graduate Texts in Mathematics \textbf{228},
Springer, New York, 2005.

\bibitem{AtkinLehner70}
A.~O.~L. Atkin and J. Lehner,
\emph{Hecke operators on $\Gamma_0(m)$},
Math.\ Ann.\ \textbf{185} (1970) 134--160.

\bibitem{Damerell71}
R.~M. Damerell,
\emph{L-functions of elliptic curves with complex multiplication, I},
Acta Arith.\ \textbf{17} (1970) 287--301;
Part~II, \textbf{19} (1971) 311--317.

\bibitem{ChowlaSelberg67}
S. Chowla and A. Selberg,
\emph{On Epstein's zeta-function},
J.~Reine Angew.\ Math.\ \textbf{227} (1967) 86--110.

\bibitem{V2_no_go}
K.~Remondi\`{e}re,
\emph{A no-go theorem for the Cabibbo angle in $S'_4$ modular flavour models
at $\tau=i$},
arXiv:[ECI v6.2 math.NT in preparation, 2026].
%% [Fix \#4: companion V2 no-go paper cited from P1 \S5.3.]

\end{thebibliography}

\end{document}
